\chapter{Analyse des besoins}\label{chap:besoins}

\section{Acteurs}

L'analyse du sujet fait apparaître trois acteurs humains et deux acteurs techniques.

\begin{table}[H]
\centering\small
\begin{tabularx}{\linewidth}{@{}l>{\raggedright\arraybackslash}X@{}}
\toprule
\textbf{Acteur} & \textbf{Rôle} \\ \midrule
Étudiant & Consulte les offres ouvertes, vérifie sa compatibilité avec chacune, candidate, suit ses candidatures, peut annuler tant que l'évaluation n'a pas commencé, consulte son stage. \\
Entreprise partenaire & Publie et gère ses offres, reçoit les candidatures, ouvre l'évaluation, accepte, refuse ou demande un complément d'information. \\
Responsable des stages & Gère les référentiels (étudiants, entreprises), consulte toutes les candidatures, suit les stages (statut, comptes rendus, évaluations), surveille les notifications et l'état des services. \\
Processus BPMN (Bonita) & Acteur technique : orchestre le traitement d'une candidature en appelant les six services. \\
Planificateur & Acteur technique : déclenche chaque jour à 8~h les rappels de date limite. \\
\bottomrule
\end{tabularx}
\caption{Acteurs du système}\label{tab:acteurs}
\end{table}

\section{Besoins fonctionnels}

\subsection{Gestion des référentiels}
\begin{itemize}
  \item Créer, consulter, modifier et supprimer un étudiant (matricule, nom, prénom, filière, niveau, coordonnées, compétences) ; changer son statut (\statut{ACTIF}, \statut{SUSPENDU}, \statut{EN\_STAGE}, \statut{DIPLOME}) selon des transitions autorisées ; vérifier son existence et son éligibilité.
  \item Créer, consulter, modifier et supprimer une entreprise (raison sociale, ICE, secteur, adresse, responsable de contact) ; consulter ses offres.
  \item Créer, modifier, clôturer et supprimer une offre (titre, description, compétences requises, filières acceptées, niveau minimum, durée, date limite, nombre de places) ; vérifier sa disponibilité et sa compatibilité avec un étudiant.
\end{itemize}

\subsection{Candidatures}
\begin{itemize}
  \item Déposer une candidature avec une lettre de motivation ; obtenir immédiatement le résultat du contrôle (enregistrée, ou rejetée avec la règle et le motif).
  \item Consulter ses candidatures, leur historique (date, auteur, commentaire de chaque changement) et les étapes du processus.
  \item Annuler une candidature lorsque les règles le permettent.
  \item Pour l'entreprise : ouvrir l'évaluation, puis accepter, refuser ou demander un complément.
\end{itemize}

\subsection{Notifications}
\begin{itemize}
  \item Confirmer le dépôt à l'étudiant ; informer l'entreprise d'une nouvelle candidature.
  \item Notifier l'étudiant d'un rejet (avec le motif), d'une acceptation, d'un refus ou d'une demande de complément.
  \item Envoyer des rappels aux étudiants compatibles qui n'ont pas encore candidaté à une offre dont la date limite approche.
\end{itemize}

\subsection{Suivi des stages}
\begin{itemize}
  \item Créer le dossier de stage d'une candidature acceptée ; consulter son statut et ses dates.
  \item Enregistrer les comptes rendus périodiques et les évaluations (note sur 20 et commentaire).
  \item Consulter l'historique des stages d'un étudiant.
\end{itemize}

\section{Besoins non fonctionnels}

\begin{table}[H]
\centering\small
\begin{tabularx}{\linewidth}{@{}l>{\raggedright\arraybackslash}X@{}}
\toprule
\textbf{Exigence} & \textbf{Réponse apportée} \\ \midrule
Autonomie des services & Un processus, un port et un schéma de base par service ; aucun accès direct aux données d'un autre service. \\
Contrats explicites & WSDL pour S1 à S3, OpenAPI pour S4 à S6, versionnés dans le dépôt ; le code est vérifié contre ces contrats par des tests. \\
Sécurité & WS-Security UsernameToken (SOAP), clé d'API par client (REST), sessions et rôles au portail, mots de passe hachés (BCrypt), masquage des coordonnées. \\
Robustesse & Délais d'attente sur tous les appels, réponse 503 explicite, sous-processus d'erreur BPMN, alerte au responsable. \\
Traçabilité & Historique de chaque changement de statut ; identifiant de corrélation propagé dans tous les appels. \\
Langue & Messages, erreurs, courriels et interface en français. \\
\bottomrule
\end{tabularx}
\caption{Besoins non fonctionnels}\label{tab:non-fonctionnels}
\end{table}

\section{Cas d'utilisation}

La figure~\ref{fig:cas-utilisation} regroupe les principaux cas d'utilisation. Le dépôt d'une candidature \emph{inclut} son traitement par le processus, qui lui-même inclut les vérifications et l'envoi des notifications.

\begin{figure}[H]
\centering
\begin{tikzpicture}[
  font=\sffamily\scriptsize,
  cas/.style={ellipse, draw=inpt, fill=inptpale, line width=0.6pt, minimum width=36mm, minimum height=7.5mm, inner sep=1pt, align=center, text width=31mm},
  acteur/.style={align=center, font=\sffamily\scriptsize},
  lien/.style={draw=encre2, line width=0.45pt},
  inclut/.style={draw=encre2, dashed, -{Stealth[length=2mm]}, line width=0.45pt},
]
\newcommand{\bonhomme}[2]{%
  \begin{scope}[shift={(#1)}]
    \draw[encre, line width=0.7pt] (0,0.42) circle (0.13);
    \draw[encre, line width=0.7pt] (0,0.29) -- (0,-0.05) (-0.2,0.18) -- (0.2,0.18) (0,-0.05) -- (-0.15,-0.3) (0,-0.05) -- (0.15,-0.3);
    \node[acteur, below] at (0,-0.32) {#2};
  \end{scope}}

% Système
\draw[filet, line width=0.8pt] (-4.2,-8.15) rectangle (4.2,5.15);
\node[font=\titres\small\bfseries, anchor=north west, text=encre2] at (-4.1,5.1) {Plateforme des stages};

% Cas étudiant
\node[cas] (offres) at (-2.05,4.2) {Consulter les offres et la compatibilité};
\node[cas] (deposer) at (-2.05,3.1) {Déposer une candidature};
\node[cas] (suivre) at (-2.05,2.0) {Suivre ses candidatures (historique, processus)};
\node[cas] (annuler) at (-2.05,0.9) {Annuler une candidature};
\node[cas] (stage) at (-2.05,-0.2) {Consulter son stage et ses courriels};

% Cas entreprise
\node[cas] (publier) at (2.05,4.2) {Publier et gérer ses offres};
\node[cas] (recevoir) at (2.05,3.1) {Consulter les candidatures reçues};
\node[cas] (evaluer) at (2.05,2.0) {Ouvrir l'évaluation};
\node[cas] (decider) at (2.05,0.9) {Décider : accepter, refuser, demander un complément};

% Processus
\node[cas, fill=white] (traiter) at (0,-1.5) {Traiter la candidature (processus BPMN)};
\node[cas, fill=white] (verifier) at (-2.05,-2.75) {Vérifier étudiant, offre, entreprise (R1 à R7)};
\node[cas, fill=white] (notifier) at (2.05,-2.75) {Notifier par courriel};
\node[cas, fill=white] (creerstage) at (2.05,-1.5) {Créer le dossier de stage (R10)};

% Cas responsable
\node[cas] (referentiels) at (-2.05,-4.15) {Gérer étudiants et entreprises};
\node[cas] (toutes) at (2.05,-4.15) {Consulter toutes les candidatures};
\node[cas] (suivistage) at (-2.05,-5.3) {Suivre les stages (comptes rendus, évaluations)};
\node[cas] (rappels) at (2.05,-5.3) {Envoyer les rappels de date limite};
\node[cas] (supervision) at (0,-6.5) {Superviser les notifications et l'état des services};
\node[cas, fill=white] (connexion) at (0,-7.55) {S'authentifier};

% Acteurs
\bonhomme{-6.0,2.2}{Étudiant}
\bonhomme{6.0,2.6}{Entreprise}
\bonhomme{-6.0,-5.0}{Responsable\\des stages}
\node[draw=encre2, font=\sffamily\scriptsize, align=center, inner sep=3pt] (bonita) at (6.0,-1.6) {«\,système\,»\\Processus\\Bonita};
\node[draw=encre2, font=\sffamily\scriptsize, align=center, inner sep=3pt] (planif) at (6.0,-5.3) {«\,système\,»\\Planificateur};

\coordinate (E) at (-5.75,2.3); \coordinate (P) at (5.75,2.7); \coordinate (R) at (-5.75,-4.9);
\foreach \c in {offres, deposer, suivre, annuler, stage} \draw[lien] (E) -- (\c.west);
\foreach \c in {publier, recevoir, evaluer, decider} \draw[lien] (P) -- (\c.east);
\foreach \c in {referentiels, suivistage, toutes, rappels, supervision} \draw[lien] (R) -- (\c);
\draw[lien] (bonita) -- (traiter.east);
\draw[lien] (planif) -- (rappels.east);

\draw[inclut] (deposer.south) .. controls (-0.8,1.2) and (-0.4,-0.6) .. node[pos=0.55, fill=white, inner sep=1pt, font=\sffamily\tiny] {«\,inclut\,»} (traiter.north);
\draw[inclut] (decider.south) .. controls (1.2,-0.2) .. (traiter.north east);
\draw[inclut] (traiter) -- (verifier);
\draw[inclut] (traiter) -- (notifier);
\draw[inclut] (traiter) -- (creerstage);
\end{tikzpicture}
\caption{Diagramme des cas d'utilisation}\label{fig:cas-utilisation}
\end{figure}

\section{Statuts manipulés}

Quatre cycles de vie structurent le métier. Leurs transitions sont contrôlées par les services eux-mêmes : une transition non prévue est refusée avec une erreur de conflit (faute \code{ConflitFault} en SOAP, code HTTP 409 en REST).

\begin{table}[H]
\centering\small
\begin{tabularx}{\linewidth}{@{}l>{\raggedright\arraybackslash}X@{}}
\toprule
\textbf{Objet} & \textbf{Statuts et transitions autorisées} \\ \midrule
Étudiant (S1) & \statut{ACTIF} $\rightarrow$ \statut{SUSPENDU}, \statut{EN\_STAGE}, \statut{DIPLOME} ; \statut{SUSPENDU} $\rightarrow$ \statut{ACTIF} ; \statut{EN\_STAGE} $\rightarrow$ \statut{ACTIF}, \statut{DIPLOME} \\
Offre (S3) & \statut{BROUILLON}, \statut{ACTIVE}, \statut{POURVUE} (plus aucune place), \statut{CLOTUREE} \\
Candidature (S4) & \statut{DEPOSEE} $\rightarrow$ \statut{EN\_COURS\_EVALUATION} $\rightarrow$ \statut{ACCEPTEE} ou \statut{REFUSEE} ; \statut{DEPOSEE} $\rightarrow$ \statut{ANNULEE} (règle R9) \\
Stage (S6) & \statut{PLANIFIE} $\rightarrow$ \statut{EN\_COURS} ou \statut{INTERROMPU} ; \statut{EN\_COURS} $\rightarrow$ \statut{TERMINE} ou \statut{INTERROMPU} \\
\bottomrule
\end{tabularx}
\caption{Cycles de vie des principaux objets}\label{tab:statuts}
\end{table}

Une \textbf{soumission} (S4) précède la candidature : c'est la demande de l'étudiant telle qu'elle entre dans le processus. Elle devient \statut{ENREGISTREE} (une candidature est alors créée), \statut{REJETEE} (avec la règle violée et le motif) ou \statut{ERREUR} (un service n'a pas répondu). Cette distinction garde une trace des dépôts refusés sans les compter comme des candidatures.
